% ECONOMICS_PROBLEM: {"id": "D71-180", "title": "Welfare with limited preference information", "jel_code": "D71", "source_id": "OP-0279"}
% !TeX program = xelatex
\documentclass[11pt,a4paper]{article}
\usepackage[margin=23mm]{geometry}
\usepackage{fontspec}
\setmainfont{Latin Modern Roman}
\usepackage{amsmath,amssymb,mathtools}
\usepackage{xcolor,enumitem,booktabs,longtable,array,xurl,hyperref}
\definecolor{muted}{HTML}{53636C}
\hypersetup{colorlinks=true,urlcolor=blue,linkcolor=blue}
\setlength{\parindent}{0pt}
\setlength{\parskip}{5pt}
\newcommand{\R}{\mathbb R}
\newcommand{\E}{\mathbb E}
\newcommand{\N}{\mathbb N}
\newcommand{\ind}{\mathbf 1}
\newcommand{\Var}{\operatorname{Var}}
\newcommand{\Cov}{\operatorname{Cov}}
\newcommand{\supp}{\operatorname{supp}}
\DeclareMathOperator*{\argmin}{arg\,min}
\DeclareMathOperator*{\argmax}{arg\,max}
\newcommand{\entrymeta}[3]{{\sffamily\small\color{muted}\textbf{\texttt{#1}}\quad|\quad #2\par #3\par}}
\newcommand{\Problemref}[1]{\texttt{#1}}
\newcommand{\JELref}[2]{\texttt{#2} (JEL \texttt{#1})}
\begin{document}
\section*{Welfare with limited preference information}
\textbf{Problem ID:} \texttt{D71-180}\quad \textbf{JEL:} \texttt{D71}\par
% BEGIN ECONOMICS STATEMENT
\label{problem:OP-0279}
{\sffamily\small\color{muted}Original subject group: Economic theory, equilibrium and social choice\par}
\label{definitions:problem:30007561}
\textbf{Audit classification:} Published research agenda. The 2021 survey’s numbered Open Problems 3–4 concern information tradeoffs. This particular normalized utility query frontier is editorial and no unspecified parameter regime is certified unsolved.
\subsection*{Definitions and precise research target}
\textbf{Model and research target.} Let \(n\) voters choose one of \(m\geq2\) alternatives. Voter \(i\) has unknown nonnegative utilities \(u_i(a)\) normalized by \(\sum_{a=1}^m u_i(a)=1\). The mechanism initially observes each voter's ordinal ranking, with a fixed rule for utility ties, and may adaptively query at most \(k\) utility values per voter. Queries return truthful exact values; strategic reporting is outside this specification. A possibly randomized algorithm \(A\) selects an alternative, and social welfare is \(W_u(a)=\sum_i u_i(a)\). Define its worst-case distortion by
\[
 D(A)=\sup_u\frac{\max_a W_u(a)}{E[W_u(A(u))]},\qquad
 D^*(n,m,k)=\inf_A D(A),
\]
where the supremum ranges over normalized utility profiles and the expectation over algorithmic randomization. Determine sharp orders and constants for \(D^*(n,m,k)\) in parameter regimes not covered by existing results, together with optimal adaptive query policies and matching information-theoretic lower bounds. Explicitly state the regime and all known upper and lower bounds before claiming a residual gap. Examine whether bounded measurement error or a restricted nonnegative valuation domain changes the query–welfare frontier.

\emph{Definitions and maintained assumptions (editorial unless a source definition is named).} Utilities lie in \(\mathcal U=\{u\in\mathbb R_+^{n\times m}:\sum_a u_i(a)=1\ \forall i\}\). Rankings must be consistent with these utilities and use a public deterministic tie order. An adaptive query policy sees prior rankings, previous exact answers and its own random seed; it may query a pair \((i,a)\) no more than \(k\) times per voter in total. Set a welfare ratio with zero denominator and positive numerator to \(+\infty\). Because \(\sum_a W_u(a)=n\), uniform random choice guarantees distortion at most \(m\); \(k=m\) gives distortion one by exact maximization, and with one voter the top ranked alternative already achieves one. These are solved baselines. Optimizing over arbitrary algorithms requires no computational limit; a polynomial-time version measures rational-answer bit length and total operations. Any claimed open residual case must fix a valuation domain, randomization convention, normalization and query range and establish a documented gap.
\subsection*{Variants and assumption changes}
\begin{enumerate}
\item \textbf{Noisy queries (editorial).} Each answer differs from its true utility by at most \(\eta\), adversarially or under a stated sampling model. Determine welfare guarantees as functions of \((m,k,\eta)\), preserving utility normalization and distinguishing additive welfare error from multiplicative distortion.
\item \textbf{Alternative setting (source agenda).} Use one-sided matching or metric social choice with precisely defined costs instead of unit-sum voting utilities. Ask how limited cardinal information improves the corresponding distortion; numerical bounds do not transfer across normalizations or objectives.
\end{enumerate}
\subsection*{References and source audit}
\begin{enumerate}
\item §5 Open Problems 3–4, printed p.9. Supports the stated research area; exact displayed specification is editorial. \url{https://arxiv.org/pdf/2103.00911}
\end{enumerate}
\begin{enumerate}\raggedright
\item\hypertarget{definitions:source:30007561:1}{} Elliot Anshelevich; Aris Filos-Ratsikas; Nisarg Shah; Alexandros A. Voudouris. 2021. \emph{Distortion in Social Choice Problems: The First 15 Years and Beyond}. §5 Generalizations of Distortion, Limited cardinal queries, Open Problems 3–4, PDF p.9.
\url{https://arxiv.org/pdf/2103.00911}.
DOI: \url{https://doi.org/10.24963/ijcai.2021/589}.
\end{enumerate}

\subsection*{Common conventions: Source mathematical conventions}

These conventions apply to each entry unless explicitly replaced.

\textbf{Models and identification.} A primitive model \(m\) specifies all probability laws, feasible choices, information, equilibrium and measurement rules used by its target. The admissible class \(\mathcal M\) is fixed independently of outcomes. If \(P_O(m)\) is its observable probability law and \(\tau(m)\) the target, its identified set at observable law \(P\) is
\[
 \mathcal I_\tau(P)=\{\tau(m):m\in\mathcal M,\ P_O(m)=P\}.
\]
Point identification means that this set is a singleton. An empty set signals incompatibility of data and assumptions. Coordinate bounds need not describe the joint set. Bounds are sharp only if every retained target is attainable or arbitrarily approachable by compatible models. Finite bounds require explicit restrictions on unobserved quantities; unrestricted confounding, hidden wealth, extrapolation or arbitrary equilibrium selection can yield uninformative sets.

\textbf{Causal interventions.} A potential outcome \(Y_i(p)\) is the outcome under a complete policy regime \(p\), including timing, financing, information and market spillovers. The target population and units are fixed before treatment. With interference, use the complete assignment vector or a justified exposure mapping; individual no-interference notation alone cannot identify an economy-wide intervention. A difference of mean potential outcomes does not require the cross-world joint distribution. Conditioning on post-treatment migration, survival, enrollment or employment changes the estimand and needs separate justification.

\textbf{Statistical statements.} An estimator, confidence set, forecast or test specifies a sampling law, model class and loss/coverage criterion. Uniform \((1-\alpha)\)-coverage means \(\inf_{m\in\mathcal M}\Pr_m\{\tau(m)\in C_n\}\geq1-\alpha\), or an explicitly asymptotic version. The whole parameter space trivially has coverage, so informativeness requires a stated width, risk or power objective. A forecasting improvement is relative to a named benchmark, information set and evaluation population.

\textbf{Equilibrium and optimization.} An equilibrium comprises feasible plans, optimality, beliefs where needed, budget constraints and all market/resource clearing. The primitives or a stated benchmark define each correspondence \(\mathcal E(p;m)\). Nonemptiness is assumed or proved, never inferred just from compact policy sets. A selected-equilibrium target uses a declared selection rule; a robust target quantifies over all permitted equilibria. Use suprema or infima unless attainment is established. An \(\varepsilon\)-optimal policy is feasible and within \(\varepsilon\) of the finite optimum value. Report an empty feasible set separately.

\textbf{Welfare and accounting.} Normative utility, interpersonal normalization, social weights, numeraire, discounting and the use of public receipts are primitives. Behavioral utility need not coincide with the welfare criterion. Household welfare includes ownership income and rents once; adding profits again to compensating variation generally double-counts them. A monetary equivalent is defined by an explicit transfer and price convention, with monotonicity and range conditions. Average effects, mechanism contributions, transfers and welfare losses are different targets. Infinite sums require convergence and dynamic feasible plans require terminal or no-Ponzi conditions.

\textbf{Source versions.} A working-paper date, revision date and journal publication year are distinguished. A DOI for a journal version is not automatically the DOI of an earlier manuscript. Locators refer to the stated version and printed pages where specified. A metadata-only check does not verify a claim about a full-text conclusion. Every entry's source subsection narrows the provenance claim accordingly.
% END ECONOMICS STATEMENT
\end{document}
